\documentclass[11pt]{article}

% ----------------------------------------------------------------------
%  Creo Advisors — Sample Marketing Collateral (ENT 105)
%  Ready-to-use copy executing the AI value-creation plan:
%  positioning, messaging, product description, and sample assets.
%  Draft by the ENT 105 team for Creo. Public-safe; figures illustrative.
%  Compile: pdflatex creo-marketing-collateral.tex  (run twice for ToC)
% ----------------------------------------------------------------------

\usepackage[utf8]{inputenc}
\usepackage[T1]{fontenc}
\usepackage{textcomp}
\usepackage[margin=1in]{geometry}
\usepackage{booktabs}
\usepackage{array}
\usepackage{enumitem}
\usepackage{titlesec}
\usepackage{xcolor}
\usepackage{tcolorbox}
\usepackage[hidelinks]{hyperref}

\definecolor{creonavy}{RGB}{17,38,73}
\definecolor{creoaccent}{RGB}{27,94,140}

\titleformat{\section}{\large\bfseries\color{creonavy}}{\thesection}{0.6em}{}
\titleformat{\subsection}{\normalsize\bfseries\color{creonavy!85!black}}{\thesubsection}{0.5em}{}
\titlespacing*{\section}{0pt}{1.5ex plus .3ex}{0.7ex}
\setlist[itemize]{leftmargin=1.3em,topsep=2pt,itemsep=2.5pt}
\setlist[enumerate]{leftmargin=1.6em,topsep=2pt,itemsep=3pt}
\setlength{\parskip}{4pt plus 1pt minus 1pt}
\setlength{\parindent}{0pt}
\setlength{\emergencystretch}{2em}

\newtcolorbox{bigidea}{colback=creonavy!6,colframe=creonavy,
  boxrule=0.8pt,arc=2pt,left=8pt,right=8pt,top=6pt,bottom=6pt}
\newtcolorbox{callout}{colback=creoaccent!8,colframe=creoaccent,
  boxrule=0.6pt,arc=2pt,left=8pt,right=8pt,top=5pt,bottom=5pt}
% Ready-to-use copy (looks like a draft asset on the page)
\newtcolorbox{copybox}[1]{colback=gray!4,colframe=gray!55!black,
  boxrule=0.5pt,arc=1pt,left=7pt,right=7pt,top=4pt,bottom=5pt,
  title={\footnotesize\bfseries\color{creonavy}#1},coltitle=creonavy,
  fonttitle=\footnotesize\bfseries,attach title to upper={\par\smallskip}}

\title{\textbf{\color{creonavy} Creo Advisors --- Sample Marketing Collateral}\\[3pt]
\large ``AI value is a map, not a multiplier'' --- positioning, messaging, and ready-to-use copy}
\author{\normalsize Draft sample assets prepared by the ENT 105 team for Creo Advisors}
\date{\normalsize 2 June 2026}

\begin{document}
\maketitle
\vspace{-1.2em}

\begin{bigidea}
\textbf{What this is.} A sample set of ready-to-use marketing assets that \emph{execute} the AI
value-creation plan --- so the strategy is concrete, on-brand, and immediately usable. Every
piece follows the one rule: \textbf{rigor over hype}. Nothing here oversells AI; the credibility
\emph{is} the restraint. All claims rest on the public record; ratios are illustrative orders of
magnitude, and self-reported market figures are attributed, not asserted.
\end{bigidea}

\tableofcontents
\vspace{0.5em}\hrule\vspace{0.7em}

% ======================================================================
\section{Positioning statement}
% ======================================================================
The formal positioning (Geoffrey Moore format) --- the single source of truth every other asset
flows from:

\begin{callout}
\textbf{For} private-equity sponsors and operating partners underwriting AI-driven value creation
--- \textbf{who} cannot reliably separate durable AI value from hype, and so risk overpaying for
deals or breaking the businesses they buy --- \textbf{Creo Advisors is} the independent, buyer-side
advisor \textbf{that} maps, function by function, what AI can actually automate and who it can
amplify, so you underwrite the real number and capture value that lasts. \textbf{Unlike} the
strategy firms, Big Four, and operating-partner shops whose AI ``diligence'' is a gateway to the AI
build they sell --- or the in-house platforms that build their own --- \textbf{Creo sells no AI
systems and no build.} That's the point: the only firm whose read on \emph{what to automate} has
nothing to sell on the other side. We tell you what AI \emph{can't} durably do, not just what it can.
\end{callout}

\textbf{Category line.} Independent, buyer-side AI due diligence for private equity --- we read the
deal; we don't build the AI.\\
\textbf{Internal short form.} \emph{Creo: the independent, buyer-side authority on what AI can
actually do in a deal --- the only one with nothing to sell on the other side.}

% ======================================================================
\section{Messaging house}
% ======================================================================
One core message on the roof; three pillars holding it up; proof under each.

\begin{callout}
\textbf{Core message.} \emph{AI value is a map, not a multiplier --- and we draw it.}
\end{callout}

\begin{center}
\renewcommand{\arraystretch}{1.3}\small
\begin{tabular}{p{2.7cm} p{4.5cm} p{5.0cm}}
\toprule
\textbf{Pillar} & \textbf{Claim} & \textbf{Proof / reason to believe} \\
\midrule
\textbf{Independent} & We're buyer-side and sell no AI systems \emph{and no build} --- our read isn't conflicted. & Every external AI-diligence provider (strategy firms, Big Four, ops shops) funnels diligence into a build they monetize; Creo has nothing to sell on the other side --- neutrality is the scarce asset in a hype market. \\
\textbf{Rigorous \& transparent} & Every call is shown, not asserted --- you can check it. & Grounded in established occupational-exposure research (Frey--Osborne; the OpenAI ``exposure $\neq$ automation'' study); reconciles to the number. \\
\textbf{Amplify, don't gut} & Durable value is a swap, not a purge. & Deploy-and-fire is why most AI programs show no return ($\sim$95\% of pilots fail; Klarna re-hired humans); we map who to amplify. \\
\bottomrule
\end{tabular}
\end{center}

\textbf{Message by buyer} (lead with the line that fits the seat):
\begin{itemize}
  \item \textbf{Deal partner (pre-deal):} ``Price on the real ratio, not the pitch --- before you overpay.''
  \item \textbf{Operating partner (returns to LPs):} ``Most AI programs show no return. We map where the durable value actually is.''
  \item \textbf{Portfolio CEO (operator):} ``Amplify your best people and collapse the busywork --- without gutting what works.''
  \item \textbf{LP / IR (amplifier):} ``Back your AI story with diligence you can defend.''
\end{itemize}

% ======================================================================
\section{Taglines, boilerplate \& elevator pitch}
% ======================================================================
\textbf{Tagline options} (lead candidate first):
\begin{itemize}
  \item \emph{AI value is a map, not a multiplier.}
  \item \emph{We tell you what AI can't do.}
  \item \emph{The only AI read with nothing to sell.}
  \item \emph{Price the real number.}
  \item \emph{Rigor over hype. Buyer-side, always.}
  \item \emph{Amplify the business --- don't gut it.}
\end{itemize}

\begin{copybox}{One-liner (for bios, intros, signatures)}
Creo Advisors is the independent, buyer-side advisor that tells private-equity sponsors what AI
can --- and can't --- durably do in a deal.
\end{copybox}

\begin{copybox}{30-second elevator pitch}
Every PE firm is chasing the AI roll-up, but most can't tell real AI value from hype --- and the
data shows it: about 95\% of enterprise AI pilots fail, and only roughly 5\% of companies capture
real value. Creo is the independent, buyer-side advisor that maps, function by function, what AI
can actually automate and who it can amplify --- so sponsors price the real number instead of the
pitch, and build value that lasts instead of gutting the business. We're the one voice that tells
you what AI \emph{can't} do.
\end{copybox}

\begin{copybox}{Boilerplate (standard --- press, LinkedIn ``About,'' deck back page)}
Creo Advisors is a Boston-based boutique advisory firm specializing in independent AI value
creation for private equity. Creo helps sponsors and their portfolio companies separate durable
AI value from hype --- mapping, function by function, what can actually be automated and who can
be amplified --- so buyers underwrite the real number and capture value that lasts. Independent
and buyer-side, Creo sells no AI systems and no build --- so, unlike every other AI-diligence
provider, it has nothing to sell on the other side of its own read: its only product is a rigorous,
transparent read on what AI can do in a deal.
\end{copybox}

% ======================================================================
\section{Product description --- the AI-Automatability Diagnostic}
% ======================================================================
The flagship offer, written as a sell sheet.

\begin{callout}
\textbf{The AI-Automatability Diagnostic} --- \emph{know what's real before you pay for it.}\\[2pt]
A fast, independent assessment of how much \emph{durable} value AI can actually create in a
target --- and whether the business can execute the change without breaking.
\end{callout}

\textbf{Who it's for.} Deal partners underwriting an AI thesis; operating partners who must show
AI returns; portfolio CEOs under pressure to ``do AI.''

\textbf{How it works --- two maps.}
\begin{enumerate}
  \item \textbf{The function map.} Every function sorted into \emph{Collapse} (pure coordination,
  near-total automation) / \emph{Amplify} (knowledge work, a few times more effective) /
  \emph{Level-up} (routine work, AI raises the floor) / \emph{Keep-human} (judgment, physical,
  relationship, regulated) --- with the real handoff ratio for each.
  \item \textbf{The reselection map.} Whether the business has operators worth amplifying, and
  management able to run the swap --- re-select for AI-fluent talent --- without breaking quality.
\end{enumerate}

\textbf{What you get.}
\begin{itemize}
  \item The function-by-function ratio map.
  \item A \textbf{reselection-readiness} score (amplifiable talent $+$ management's capacity to execute).
  \item The ``real vs.\ hyped value'' gap --- translated into what the deal is actually worth.
  \item The amplify / collapse / keep priorities --- and the functions (or deals) to walk away from.
\end{itemize}

\textbf{The engagement ladder.}
\begin{center}
\renewcommand{\arraystretch}{1.2}\small
\begin{tabular}{p{2.6cm} p{9.6cm}}
\toprule
\textbf{Free} & Complimentary mini-Diagnostic on one target --- the proof, and the lead magnet. \\
\textbf{Productized} & Fixed-fee full Diagnostic on a deal or portfolio company. \\
\textbf{Engagement} & The value-creation / amplification roadmap (Creo orchestrates; vendors build). \\
\textbf{Recurring} & The trusted AI read across the sponsor's portfolio. \\
\bottomrule
\end{tabular}
\end{center}

\textbf{Why Creo.} Independent (we don't sell the build) $\cdot$ rigorous and transparent
(grounded in established automation research; shown, not asserted) $\cdot$ focused on durable
value (amplify, don't gut). \emph{And the loop closes:} post-deal, Creo measures independently
whether the projected value was actually created.

% ======================================================================
\section{The AI Reality Report --- lead-magnet copy}
% ======================================================================
\begin{copybox}{Gated-asset landing copy}
\textbf{The AI Reality Report} --- what actually works, and what actually doesn't.\\[3pt]
Creo's annual, evidence-based read on where AI is creating real value in private equity --- the
function ratios that hold up, the failure rates nobody quotes, and the deals that blow up. No
hype, no vendor spin --- the independent, buyer-side view.\\[3pt]
\textbf{[ Download the report $\rightarrow$ ]} \quad\emph{(name + work email)}
\end{copybox}

% ======================================================================
\section{Website hero copy}
% ======================================================================
\begin{copybox}{Homepage hero}
\textbf{\large AI value is a map, not a multiplier.}\\[4pt]
Creo is the independent, buyer-side advisor that tells PE sponsors what AI can --- and can't ---
durably do in a deal. Price the real number. Capture value that lasts.\\[4pt]
\textbf{[ Get a free AI-Automatability Diagnostic ]} \quad \emph{[ Read the AI Reality Report ]}
\end{copybox}

\textbf{Three value props} (the band beneath the hero):
\begin{itemize}
  \item \textbf{Independent.} We sell no AI systems --- so our read isn't conflicted.
  \item \textbf{Rigorous.} Grounded in real automation research, reconciled to the number, shown not asserted.
  \item \textbf{Durable.} We amplify the people who run the business --- we don't gut it.
\end{itemize}

% ======================================================================
\section{Sample LinkedIn posts (founder-led)}
% ======================================================================
Contrarian, evidence-based, in a founder's voice --- attention \emph{and} credibility.

\begin{copybox}{Post 1 --- ``Map, not multiplier''}
Every deck I see prices AI at one number. ``AI makes this 10x more productive'' --- one multiple,
applied to the whole company.\\[3pt]
That's not how it works. A business is a portfolio of functions, and AI hits each one
completely differently:\\[2pt]
$\bullet$ Pure coordination --- dispatching, matching, routing --- can nearly disappear. (Uber
did this to an entire industry.)\\
$\bullet$ Knowledge work amplifies a few times over: a great analyst with AI does the work of
several --- but you still need the analyst.\\
$\bullet$ Judgment, relationships, regulated work, anything physical --- barely moves.\\[3pt]
Price the average and you either overpay for the hype or miss the rare function that actually
collapses. The value --- and the risk --- live in the mix.\\[3pt]
AI value is a map, not a multiplier. Most buyers are still pricing the multiplier.
\end{copybox}

\begin{copybox}{Post 2 --- ``The rehire is the tell'' (public-record only)}
The headlines say AI is replacing workers. Watch what the companies actually do.\\[3pt]
Microsoft cut more than 15,000 roles last year. It also made AI use ``no longer optional,'' tied
it to performance reviews, and shielded its own AI teams from the hiring freeze. That's not ``the
machine took the jobs.'' That's a workforce being re-sorted for who can amplify with AI.\\[3pt]
Here's the tell that separates real automation from hype: the rehire fraction. When a company
sheds thousands and then quietly rehires a sizable share, AI didn't replace those people --- it
raised the bar on who gets to stay. A big rehire fraction means amplification, not the
1000-to-1 the hype promises.\\[3pt]
For anyone underwriting an AI roll-up, that's the whole question: are you buying replacement, or
reselection?
\end{copybox}

\begin{copybox}{Post 3 --- ``95\% fail''}
About 95\% of corporate AI pilots fail to deliver. That's not my number --- it's the research.\\[3pt]
So why does the handful that succeed succeed? It's almost never the model. The winners did three
unglamorous things:\\[2pt]
1. Found the functions where AI actually compounds --- and ignored the ones where it doesn't.\\
2. Amplified their best people instead of replacing their cheapest.\\
3. Measured the result honestly, and killed what didn't work.\\[3pt]
The losers bought a platform, fired a department, and waited for margin that never came.\\[3pt]
If you're paying an AI-roll-up multiple, you're betting you'll be in the 5\%. The only way to
know is to look --- function by function --- before you sign.
\end{copybox}

\begin{copybox}{Post 4 --- ``The model just went private'' (timely; public-record only)}
This week a16z funded a private-sector DOGE. Ex-DOGE staffers launched a firm called \textbf{Special}
with one line: buy a company, automate the overhead, keep the margin. Their first target? A Texas
eldercare business with hundreds of nurses.\\[3pt]
Hold that up against the record. DOGE says it saved \$160B. One nonpartisan analysis (the
Partnership for Public Service) estimates its actions will \emph{cost} about \$135B this year once
you count paid leave, re-hiring, and lost productivity. The cut booked as savings --- and came back
as cost.\\[3pt]
Now run that logic on nurses. Eldercare is judgment, physical, relationship, regulated work --- the
most \emph{keep-human} function there is. The durable AI value in that business isn't the care; it's
the back office: scheduling, intake, billing. Automate \emph{those} and margin expands. Cut the
nurses and you haven't cut cost --- you've cut capacity, and it returns as turnover, quality, and risk.\\[3pt]
``Buy a company, automate the overhead'' is the right instinct. \emph{Which} overhead is the whole
question --- and it's a map, not a multiplier. We'll be watching this one.
\end{copybox}

% ======================================================================
\section{Sample nurture email}
% ======================================================================
\begin{copybox}{Email --- sent after an AI Reality Report download}
\textbf{Subject:} The 5\% --- what separates the AI programs that actually work\\[4pt]
Hi \{First name\},\\[3pt]
Thanks for downloading the AI Reality Report. One finding tends to stop people: roughly 95\% of
enterprise AI pilots fail to deliver a return --- and the 5\% that work rarely have better
technology. They have a better \emph{map}.\\[3pt]
That's the idea behind our AI-Automatability Diagnostic: an independent, buyer-side read of where
AI can actually create durable value in a specific business --- function by function --- and
whether the team can execute it without breaking what works.\\[3pt]
If you're underwriting an AI thesis right now, we'll run a \textbf{complimentary mini-Diagnostic}
on one target. No build pitch --- we don't sell AI systems. Just an honest read on what's real.\\[3pt]
Worth a 20-minute call?\\[3pt]
Rich Vitaro --- Creo Advisors\\
\textbf{[ Book a call ]}
\end{copybox}

\vspace{4pt}\hrule\vspace{4pt}
{\footnotesize\itshape Draft sample collateral prepared by the ENT 105 team for Creo Advisors.
Copy is illustrative and for management review; names, figures, and claims to be finalized with
Creo. All assets are public-record-only and independent --- the brand promise is rigor over hype,
so the copy never claims more about AI than the evidence supports.}

\end{document}
